% !TEX root = ../main.tex
\chapter{\ifenglish Introduction\else บทนำ\fi}

\section{\ifenglish Project rationale\else ที่มาและความสำคัญของโครงงาน\fi}
การลงคะแนนเสียงเลือกตั้งเป็นกลไกพื้นฐานที่สำคัญยิ่งในระบอบประชาธิปไตย เพื่อสะท้อนเจตจำนงของประชาชนในการเลือกตัวแทนเข้าไปบริหารประเทศ อย่างไรก็ตาม กระบวนการเลือกตั้งแบบดั้งเดิมที่ใช้บัตรกระดาษมักประสบปัญหาข้อผิดพลาดหลายประการ เช่น ปัญหาบัตรเสีย ข้อผิดพลาดในการตรวจนับคะแนน ความล่าช้าในการรายงานผล ตลอดจนปัญหาความโปร่งใสและงบประมาณในการจัดการเลือกตั้ง~\cite{ect_evoting_1, ect_evoting_2}

เพื่อแก้ไขปัญหาดังกล่าว หลายประเทศรวมถึงประเทศไทยได้เริ่มศึกษาและนำระบบการลงคะแนนเสียงอิเล็กทรอนิกส์ (Electronic Voting หรือ e-Voting) และการลงคะแนนเสียงออนไลน์ผ่านเครือข่ายอินเทอร์เน็ต (Online Internet Voting) มาปรับใช้ เพื่อเพิ่มความสะดวก รวดเร็ว และลดต้นทุนการดำเนินงาน~\cite{elect_in_th, e2e_voting} ทว่าระบบการลงคะแนนเสียงออนไลน์ในปัจจุบันส่วนใหญ่ยังคงพึ่งพาการเข้ารหัสแบบกุญแจอสมมาตร (Asymmetric Public-Key Cryptography) เช่น RSA~\cite{rivest_rsa} และการเข้ารหัสเส้นโค้งวงรี (Elliptic Curve Cryptography: ECC) เป็นโครงสร้างพื้นฐานด้านความมั่นคงปลอดภัย ซึ่งความปลอดภัยของระบบเหล่านี้ตั้งอยู่บนสมมติฐานเชิงคำนวณ (Computational Security) ว่าคอมพิวเตอร์แบบดั้งเดิมไม่สามารถแยกตัวประกอบของจำนวนเต็มขนาดใหญ่หรือคำนวณลอการิทึมไม่ต่อเนื่องได้ภายในเวลาที่เหมาะสม

อย่างไรก็ดี ความก้าวหน้าอย่างรวดเร็วของเทคโนโลยีคอมพิวเตอร์ควอนตัม (Quantum Computing) กำลังก่อให้เกิดภัยคุกคามร้ายแรงต่อระบบความปลอดภัยดังกล่าว เนื่องจากการนำอัลกอริทึมของชอร์ (Shor's Algorithm)~\cite{shor_1994, nielsen_chuang} มาประมวลผลบนเครื่องคอมพิวเตอร์ควอนตัม จะสามารถแก้ปัญหาการแยกตัวประกอบจำนวนเฉพาะและคำนวณลอการิทึมไม่ต่อเนื่องได้ภายในเวลาพหุนาม (Polynomial Time) ซึ่งส่งผลให้การเข้ารหัสกุญแจสาธารณะเกือบทั้งหมดที่ใช้อยู่ในปัจจุบันถูกถอดรหัสได้อย่างสมบูรณ์ แม้ว่าจะมีการพัฒนาระบบการเข้ารหัสยุคหลังควอนตัม (Post-Quantum Cryptography: PQC)~\cite{nist_pqc} ขึ้นมาทดแทน แต่วิธีการดังกล่าวยังคงเป็นความปลอดภัยเชิงคำนวณที่อาจเสื่อมถอยลงได้ตามพลังการคำนวณในอนาคต

เพื่อสร้างระบบการสื่อสารที่มีความปลอดภัยอย่างแท้จริงและคงทนถาวร โครงงานนี้จึงนำเสนอการประยุกต์ใช้ **การกระจายกุญแจแบบควอนตัม (Quantum Key Distribution: QKD)** ภายใต้โพรโทคอล BB84~\cite{bb84, bennett_1992_experimental, gisin_qkd_review} ผสานร่วมกับ **การเข้ารหัสแบบวันไทม์แพด (One-Time Pad: OTP)**~\cite{vernam_cipher, shannon_secure} ในการรักษาความปลอดภัยของระบบการส่งบัตรลงคะแนนเสียง โดยอาศัยกฎพื้นฐานทางกลศาสตร์ควอนตัม ได้แก่ ทฤษฎีบทการห้ามคัดลอกสถานะควอนตัม (No-Cloning Theorem)~\cite{wootters_zurek_nocloning} และการรบกวนจากการวัดค่า (Quantum Measurement Disturbance) ซึ่งรับประกันว่า หากมีผู้ดักฟัง (Eve) พยายามดักจับหรือคัดลอกโฟตอนในช่องสัญญาณควอนตัม สถานะควอนตัมจะเกิดการยุบตัวและสร้างอัตราความผิดพลาด (Quantum Bit Error Rate: QBER) ที่สูงขึ้นอย่างเห็นได้ชัด ทำให้ผู้ส่ง (Alice) และผู้รับ (Bob) สามารถตรวจจับการดักฟังและยกเลิกเซสชันได้ทันทีก่อนการส่งข้อมูลจริง ซึ่งแนวคิดดังกล่าวเคยได้รับการพิสูจน์เชิงประจักษ์ในการส่งผลคะแนนเลือกตั้งรัฐสภาแห่งรัฐเจนีวา ประเทศสวิตเซอร์แลนด์ ในปี ค.ศ. 2007~\cite{idquantique_geneva_2007} รวมถึงงานวิจัยด้านโปรโตคอลการลงคะแนนนิรนามเชิงควอนตัม~\cite{vaccaro_quantum_voting, horoshko_quantum_voting}

โครงงานนี้จึงมุ่งเน้นการออกแบบและพัฒนาระบบจำลองการลงคะแนนเสียงออนไลน์เชิงการศึกษา (Educational Interactive Web Simulator) ที่บูรณาการขั้นตอนการแลกเปลี่ยนกุญแจควอนตัม BB84 การสละบิตตัวอย่างเพื่อประเมินค่า QBER การขยายความลับด้วยเมทริกซ์โทพลิตซ์ (Toeplitz Privacy Amplification) และการเข้ารหัสบัตรลงคะแนนแบบ One-Time Pad เพื่อให้เห็นภาพการทำงานจริงทั้งบนโปรแกรมจำลองควอนตัมและคอมพิวเตอร์ควอนตัมจริงของ IBM Quantum ผ่านระบบเว็บแอปพลิเคชันที่มีความโปร่งใสและตรวจสอบย้อนหลังได้

\section{\ifenglish Objectives\else วัตถุประสงค์ของโครงงาน\fi}
\begin{enumerate}
    \item เพื่อศึกษา วิเคราะห์ และประยุกต์ใช้โพรโทคอลการกระจายกุญแจแบบควอนตัม BB84 ร่วมกับการเข้ารหัส One-Time Pad ในการรักษาความปลอดภัยของระบบลงคะแนนเสียงออนไลน์
    \item เพื่อพัฒนาระบบเว็บแอปพลิเคชันจำลองการลงคะแนนเสียง ที่แสดงขั้นตอนการทำงานของวงจรควอนตัม การส่งผ่านข้อมูล และการประเมินความปลอดภัยของกุญแจลับก่อนนำไปใช้เข้ารหัส
    \item เพื่อจำลองและเปรียบเทียบพฤติกรรมความปลอดภัยของระบบในสภาวะต่าง ๆ ได้แก่ ช่องสัญญาณอุดมคติ (Ideal Simulator), ช่องสัญญาณที่มีสัญญาณรบกวนจากฮาร์ดแวร์ควอนตัมจริง (IBM Quantum Cloud Processor), และสภาวะที่มีการโจมตีแบบดักวัดและส่งต่อโฟตอนของผู้ดักฟัง (Active Eve Intercept-and-Resend)
    \item เพื่อศึกษาและหาค่าพารามิเตอร์ที่เหมาะสมในการสร้างกุญแจ ได้แก่ ค่าเผื่อความปลอดภัย ($\delta$ Safety Margin) และพารามิเตอร์ความปลอดภัยในการขยายความลับ (Security Parameter $s$) เพื่อรับประกันความน่าเชื่อถือและความลับสมบูรณ์แบบ
\end{enumerate}

\section{\ifenglish Project scope\else ขอบเขตของโครงงาน\fi}

\subsection{\ifenglish Hardware scope\else ขอบเขตด้านฮาร์ดแวร์\fi}
\begin{enumerate}
    \item เนื่องจากคอมพิวเตอร์ควอนตัมยังไม่มีการใช้งานทั่วไปในระดับอุปกรณ์ส่วนบุคคล โครงงานนี้จึงใช้เครื่องคอมพิวเตอร์แม่ข่ายและคอมพิวเตอร์ลูกข่ายทั่วไปที่เชื่อมต่อผ่านระบบเครือข่ายอินเทอร์เน็ต
    \item เชื่อมต่อและรันวงจรควอนตัมบนฮาร์ดแวร์หน่วยประมวลผลควอนตัมจริง (Quantum Processing Unit: QPU) ของ IBM Quantum ผ่านระบบคลาวด์ ภายใต้ข้อจำกัดฮาร์ดแวร์รีจิสเตอร์ขนาดไม่เกิน 40 คิวบิตต่อวงจร
\end{enumerate}

\subsection{\ifenglish Software scope\else ขอบเขตด้านซอฟต์แวร์\fi}
\begin{enumerate}
    \item พัฒนาระบบประมวลผลวงจรควอนตัมด้วยภาษา Python โดยใช้ไลบรารี \texttt{qiskit-aer} สำหรับจำลองวงจรในระดับเครื่อง และ \texttt{qiskit-ibm-runtime}~\cite{qiskit2024} สำหรับส่งงานประมวลผลไปยัง IBM Quantum
    \item จำลองกระบวนการประมวลผลหลังการวัดค่า (Post-Processing Pipeline) ประกอบด้วย การคัดกรองฐานวัด (Basis Sifting), การสละบิตตัวอย่างเพื่อประเมินค่า QBER เทียบกับเกณฑ์ความปลอดภัยร้อยละ 11.0, และการขยายความลับด้วยการคูณเมทริกซ์โทพลิตซ์บนสนามจำกัด $\text{GF}(2)$ (Toeplitz Matrix Universal Hashing)
    \item พัฒนาระบบส่วนติดต่อผู้ใช้งาน (Frontend) ในรูปแบบ Single Page Application (SPA) ด้วยภาษา TypeScript ร่วมกับไลบรารี React และ Vite เพื่อแสดงผลการทำงานแบบโต้ตอบ (Interactive UI) 3 มุมมอง ได้แก่ มุมมองผู้ลงคะแนน (Voter View), มุมมองตามความเป็นจริง (Practical View), และมุมมองทั้งระบบ (System God View)
    \item พัฒนาระบบบริการส่วนหลัง (Backend RESTful API) ด้วย FastAPI และจัดเก็บข้อมูลเซสชันการลงคะแนนในฐานข้อมูล MongoDB
\end{enumerate}

\section{\ifenglish Expected outcomes\else ประโยชน์ที่คาดว่าจะได้รับ\fi}
\begin{enumerate}
    \item ได้แบบจำลองระบบลงคะแนนเสียงออนไลน์ที่สามารถพิสูจน์แนวคิด (Proof of Concept) การรักษาความปลอดภัยเชิงทฤษฎีสารสนเทศด้วยการกระจายกุญแจควอนตัมและ One-Time Pad
    \item ได้ระบบเว็บแอปพลิเคชันเชิงการศึกษาที่สามารถใช้เป็นสื่อการเรียนรู้กระบวนการทำงานของโพรโทคอล BB84 ปรากฏการณ์การยุบตัวของสถานะควอนตัม และการตรวจจับผู้ดักฟัง
    \item ได้ผลการวิเคราะห์และข้อแนะนำเชิงวิศวกรรมเกี่ยวกับขนาดการส่งควอนตัมบิต ($N$) ค่าเผื่อความปลอดภัย ($\delta$) และการเชื่อมต่อกุญแจหลายรอบ (Multi-Session Key Concatenation) สำหรับการนำ QKD ไปประยุกต์ใช้จริงภายใต้ข้อจำกัดของฮาร์ดแวร์ควอนตัมในปัจจุบัน
\end{enumerate}

\section{\ifenglish Technology and tools\else เทคโนโลยีและเครื่องมือที่ใช้\fi}

\subsection{\ifenglish Hardware technology\else เทคโนโลยีด้านฮาร์ดแวร์\fi}
\begin{enumerate}
    \item เครื่องคอมพิวเตอร์ส่วนบุคคลสำหรับการประมวลผลแม่ข่ายและทดสอบระบบ
    \item ระบบประมวลผลคอมพิวเตอร์ควอนตัมจริงของ IBM Quantum Cloud (IBM Quantum Eagle / Heron Architecture)
\end{enumerate}

\subsection{\ifenglish Software technology\else เทคโนโลยีด้านซอฟต์แวร์\fi}
\begin{enumerate}
    \item \textbf{ภาษาและสภาพแวดล้อมการพัฒนา:} Python (เวอร์ชัน 3.11 ขึ้นไป), TypeScript, Node.js
    \item \textbf{เครื่องมือทางควอนตัมคอมพิวติ้ง:} Qiskit SDK, \texttt{qiskit-aer}, \texttt{qiskit-ibm-runtime}~\cite{qiskit2024}
    \item \textbf{ระบบบริการส่วนหลังและฐานข้อมูล:} FastAPI, Uvicorn, MongoDB, Beanie ODM, Pydantic V2
    \item \textbf{ระบบส่วนติดต่อผู้ใช้งาน:} React 18, Vite, React Router, CSS Modules
    \item \textbf{การออกแบบและควบคุมเวอร์ชัน:} Figma (การออกแบบ UI/UX), Git และ GitHub
\end{enumerate}

\section{\ifenglish Budget plan\else แผนการใช้จ่ายงบประมาณ\fi}
โครงงานนี้เน้นการวิจัยและพัฒนาบนระบบซอฟต์แวร์จำลองและบริการระบบคลาวด์ระดับใช้งานฟรี (Free Open-Plan / Academic Access) ของ IBM Quantum Platform จึงไม่มีค่าใช้จ่ายด้านอุปกรณ์ฮาร์ดแวร์ควอนตัมโดยตรง ทั้งนี้ อาจมีค่าใช้จ่ายเบ็ดเตล็ดสำหรับการจัดทำเอกสารและรูปเล่มรายงานประมาณ 1,500 บาท

\section{\ifenglish Project plan\else แผนการดำเนินงาน\fi}
แผนการดำเนินงานของโครงงานแบ่งออกเป็น 4 ระยะ ในช่วงปีการศึกษา 2568--2569 ดังแสดงในตารางแผนภูมิแกนต์:

\begin{plan}{6}{2025}{3}{2026}
    \planitem{6}{2025}{8}{2025}{ศึกษาค้นคว้าทฤษฎีควอนตัม BB84, One-Time Pad และงานวิจัยที่เกี่ยวข้อง}
    \planitem{8}{2025}{10}{2025}{ออกแบบสถาปัตยกรรมระบบ โครงสร้างข้อมูล และส่วนติดต่อผู้ใช้งานบน Figma}
    \planitem{10}{2025}{1}{2026}{พัฒนาระบบจำลองวงจรควอนตัม อัลกอริทึมหลังการวัดค่า และเว็บแอปพลิเคชัน}
    \planitem{1}{2026}{3}{2026}{ทดสอบระบบ ประเมินประสิทธิภาพ วิเคราะห์ผลการทดลอง และจัดทำรายงานฉบับสมบูรณ์}
\end{plan}

\section{\ifenglish Roles and responsibilities\else บทบาทและความรับผิดชอบของคณะผู้จัดทำ\fi}
\begin{enumerate}
    \item \textbf{นายภคิน นิรมิตสีมา (รหัสนักศึกษา 660610782):} รับผิดชอบการพัฒนาระบบส่วนติดต่อผู้ใช้งาน (Frontend Web Application) ด้วยภาษา TypeScript และ React, การสร้างกลไกการนำเสนอข้อมูลแบบโต้ตอบ (Interactive Stepper), การจัดการสถานะของแอปพลิเคชัน และระบบแสดงผลหน้าบันทึกการตรวจสอบย้อนหลัง
    \item \textbf{นายยุทธกานต์ สาจุ้ย (รหัสนักศึกษา 660610787):} รับผิดชอบการพัฒนาระบบประมวลผลฝั่งเซิร์ฟเวอร์ (Backend RESTful API) ด้วย FastAPI, การออกแบบโครงสร้างฐานข้อมูล MongoDB สำหรับจัดการเซสชันการเลือกตั้งและบันทึกประวัติการสื่อสารควอนตัม (Quantum Audit Log), และการเชื่อมต่อกระบวนการถอดรหัสผลคะแนน One-Time Pad
    \item \textbf{นายอาทินันท์ แจ้งสว่าง (รหัสนักศึกษา 660610807):} รับผิดชอบการออกแบบประสบการณ์ผู้ใช้และหน้าตาส่วนติดต่อผู้ใช้งาน (UI/UX Design) บนโปรแกรม Figma, การพิสูจน์ความปลอดภัยของกุญแจที่เชื่อมต่อจากหลายรอบ, การจำลองพฤติกรรมการโจมตีของผู้ดักฟัง (Eve Intercept-and-Resend) บนหน้าเว็บ, และการจัดทำเอกสารคู่มือการใช้งานระบบ
    \item \textbf{นายไพสิฐ เลิศอนันต์พิพัฒน์ (รหัสนักศึกษา 660612155):} รับผิดชอบการพัฒนาระบบจำลองวงจรควอนตัมด้วย Qiskit, การเขียนสคริปต์เชื่อมต่อกับ IBM Quantum Runtime, การพัฒนาโมดูลการประมวลผลหลังการวัดค่า (Toeplitz Privacy Amplification), รวมถึงการออกแบบชุดทดสอบและรวบรวมข้อมูลผลการทดลองเชิงสถิติ

\end{enumerate}

\section{\ifenglish%
Impacts of this project on society, health, safety, legal, and cultural issues
\else%
ผลกระทบด้านสังคม สุขภาพ ความปลอดภัย กฎหมาย และวัฒนธรรม
\fi}

\begin{enumerate}
    \item \textbf{ผลกระทบด้านสังคมและการศึกษา (Social and Educational Impact):} ระบบจำลองนี้เป็นเครื่องมือเชิงการศึกษาที่มีประโยชน์อย่างยิ่งในการช่วยให้นักศึกษา นักวิจัย และบุคคลทั่วไปสามารถทำความเข้าใจหลักการทำงานของฟิสิกส์ควอนตัมและการเข้ารหัสลับเชิงควอนตัมที่ซับซ้อนได้อย่างเป็นรูปธรรมผ่านการจำลองเชิงภาพ (Interactive Visualization) ซึ่งช่วยส่งเสริมการสร้างบุคลากรที่มีความรู้ความเข้าใจในเทคโนโลยีควอนตัมของประเทศ
    \item \textbf{ผลกระทบด้านความมั่นคงปลอดภัย (Security and Privacy Impact):} โครงงานนี้นำเสนอแนวทางการป้องกันความปลอดภัยทางไซเบอร์ที่สามารถต้านทานการโจมตีจากคอมพิวเตอร์ควอนตัมในอนาคตได้อย่างเด็ดขาด (Quantum-Resistant) การรับประกันความปลอดภัยเชิงทฤษฎีสารสนเทศ (Information-Theoretic Security) จะช่วยสร้างความเชื่อมั่นอย่างสูงสุดให้แก่ระบบการลงคะแนนเสียงว่าความลับของบัตรลงคะแนนจะไม่สามารถถูกเจาะทำลายย้อนหลังได้ตลอดไป
    \item \textbf{ผลกระทบด้านกฎหมายและความโปร่งใส (Legal and Transparency Impact):} การมีระบบบันทึกประวัติการสื่อสารควอนตัม (Quantum Audit Log) เปิดโอกาสให้หน่วยงานกำกับดูแลการเลือกตั้งและประชาชนสามารถร่วมกันตรวจสอบความโปร่งใสของกระบวนการส่งข้อมูลได้ โดยไม่ละเมิดสิทธิ์ในการรักษาความลับของผู้ลงคะแนนเสียง ซึ่งสอดคล้องกับหลักธรรมาภิบาลและกฎหมายการเลือกตั้งสากล
\end{enumerate}
